\chapter{The Examples Catalog and Verification Strategy}
\label{ch:examples-catalog}

The repository documents its runnable and research-oriented examples
(\texttt{examples/}) as a guided catalog rather than a replacement for the
source files themselves: each page explains scientific/computational
purpose, the DAGonStar mechanisms exercised, prerequisites, execution model,
and expected verification strategy. This chapter surveys that catalog as a
reference map for readers extending this book's lessons toward their own
research problems.

\section{The Course Examples}

The course ships one folder per lecture under \texttt{examples/}, fifty
scripts in all. Every script uses the shared dataset and helpers of
\texttt{examples/common/}, runs in a fresh temporary scratch directory,
ends with its own assertions, and has its expected output recorded in
\texttt{output/<script>.txt}. Services that are not on a laptop (Docker,
Slurm, Kubernetes, Nomad, SSH hosts, clouds, live LLM, FaaS and IoT
providers, DAGonFS mounts) are replaced by mocks, local stand-ins or portable
emulation, and the scripts say so in their output.

\begin{lstlisting}[style=bash]
python tools/run_examples.py            # run all and compare with the recorded outputs
python tools/run_examples.py 07 21      # only lectures 07 and 21
python tools/run_examples.py --record   # rewrite output/*.txt and the READMEs
\end{lstlisting}

{\small
\begin{longtable}{@{}p{0.43\linewidth}p{0.53\linewidth}@{}}
\caption{The course examples.}\label{tab:course-examples}\\
\toprule
\textbf{Script} & \textbf{What it shows} \\
\midrule
\endfirsthead
\toprule
\textbf{Script} & \textbf{What it shows} \\
\midrule
\endhead
\texttt{01-hello/hello\_workflow.py} & the smallest useful workflow, one task that fetches the data \\
\texttt{02-script-vs-dag/as\_dag.py} & the same analysis as a script and as a DAG of tasks \\
\texttt{03-architecture/factory.py} & the task factory, the task classes and the enumerations \\
\texttt{04-environment/check\_environment.py} & check the installation without any configuration file \\
\texttt{05-configuration/configuration.py} & the same configuration as a dagon.ini file and as a dictionary \\
\texttt{06-first-task/first\_task.py} & one local batch task, its status and its working directory \\
\texttt{07-dag/fan\_out\_fan\_in.py} & a fan-out/fan-in DAG built with explicit control dependencies \\
\texttt{08-dataflow/workflow\_references.py} & a data dependency expressed with a workflow:// reference \\
\texttt{09-cycles/cycles.py} & a cycle is rejected before anything runs, then repaired \\
\texttt{10-scratch/inspect\_scratch.py} & what DAGonStar leaves in a task's working directory \\
\texttt{11-inspect/inspect\_workflow.py} & serialise a workflow to JSON, reload it, time it, observe it \\
\texttt{12-station-pipeline/station\_pipeline.py} & the Gulf of Naples station pipeline, end to end \\
\texttt{13-staging/staging\_modes.py} & COPY, LINK and DONTMOVE: how an input reaches the consumer \\
\texttt{14-garbage/garbage\_collection.py} & reference counting and the scratch garbage collector \\
\texttt{15-checkpoint/checkpoint\_resume.py} & record a checkpoint, then resume without re-running tasks \\
\texttt{16-failures/fail\_and\_resume.py} & a failing task, what it stops, and the resumed run \\
\texttt{17-async/async\_events.py} & launch in the background and observe the lifecycle events \\
\texttt{18-native/native\_task.py} & a native task runs an importable Python function \\
\texttt{19-native-bindings/native\_bindings.py} & scalar inputs, environment variables and a failing function \\
\texttt{20-sweep/parameter\_sweep.py} & generate a workflow from parameters: one summary per station \\
\texttt{21-ensemble/ensemble.py} & an ensemble of perturbed members and a fan-in statistic \\
\texttt{22-meta/meta\_workflow.py} & two workflows composed by DAG\_TPS with a cross-workflow reference \\
\texttt{23-dry-portable/dry\_and\_portable.py} & a dry run plans without executing; portable emulation runs\ldots \\
\texttt{24-resumable-ensemble/resumable\_ensemble.py} & a resumable ensemble, and the limit of name-based checkpoints \\
\texttt{25-web-get/web\_get.py} & a web task performs one declarative HTTP request \\
\texttt{26-web-upload/web\_upload.py} & upload a workflow artifact, authenticate from the environment, and\ldots \\
\texttt{27-llm/llm\_task.py} & an LLM task turns workflow data into a prompt and saves the reply \\
\texttt{28-llm-audit/llm\_audit.py} & what an LLM task records, what it never records, and a\ldots \\
\texttt{29-faas/faas\_mock.py} & a FaaS task invokes an already-deployed function (mock provider) \\
\texttt{30-faas-http/faas\_http.py} & the same FaaS contract against an HTTP function, and an unknown\ldots \\
\texttt{31-iot-observe/iot\_observe.py} & a bounded IoT observation and its checkpoint \\
\texttt{32-iot-actuate/iot\_actuate.py} & desired-state configuration, actuation, and an uncertain actuation\ldots \\
\texttt{33-agentic-contract/agentic\_contract.py} & the agentic contract, its dagon.ini profile, and the configurations\ldots \\
\texttt{34-agentic-guardrails/agentic\_guardrails.py} & declared tools, guardrails checked outside the agent, and an\ldots \\
\texttt{35-agentic-evaluation/agentic\_evaluation.py} & independent result fit, mandatory gates, publication only when\ldots \\
\texttt{36-continuum/continuum.py} & one graph across the continuum: an IoT observation at the edge, a\ldots \\
\texttt{37-docker/docker\_task.py} & a Docker task wraps the command in a container; the client is\ldots \\
\texttt{38-apptainer/apptainer\_task.py} & Apptainer image preparation and execution, subprocess mocked \\
\texttt{39-orchestrators/kubernetes\_nomad.py} & a pod on Kubernetes and a job on Nomad, both APIs mocked \\
\texttt{40-remote/ssh\_and\_cloud.py} & remote tasks over SSH and cloud tasks, with no network \\
\texttt{41-slurm/slurm\_tasks.py} & Slurm submissions: the generated sbatch command and two jobs run\ldots \\
\texttt{42-beeond/beeond.py} & BeeOND job-local storage: validated configuration, the sbatch\ldots \\
\texttt{43-prestage/prestage.py} & stage inputs before sbatch, verify them, and never submit when\ldots \\
\texttt{44-movers/data\_movers.py} & which data mover the stager chooses, and the commands it generates\ldots \\
\texttt{45-fair/fair.py} & FAIR metadata is opt-in, declared, licensed and credential-free \\
\texttt{46-cwl/cwl\_export.py} & a deterministic CWL v1.2 export: file edges, ordering edges, no\ldots \\
\texttt{47-operations/operations.py} & operating DAGonStar: service registration (mocked), JSON round\ldots \\
\texttt{48-capstone/capstone.py} & every registered TaskType in one checkpointed, FAIR-recorded graph\ldots \\
\texttt{A1-dagonfs/block\_layout.py} & the DAGonFS block layout, computed from (S, P) alone \\
\texttt{A1-dagonfs/dagonfs\_workflow.py} & a DAGonStar workflow whose scratch lives on a DAGonFS volume \\
\texttt{A2-volumes/local\_directory.py} & an ad-hoc adapter; workflow- and task-scope volumes \\
\texttt{A2-volumes/volume\_task.py} & volume tasks: create, share, stream, destroy, reference counting \\
\texttt{A3-meta/meta\_task.py} & meta-tasks: a command-line engine, a CWL tool, an emulated Nextflow \\
\texttt{A4-dagonweb/dagonweb\_document.py} & a DAGonWeb document loaded, run and resumed by DAGonStar \\
\bottomrule
\end{longtable}
}

The upstream repository (\url{https://github.com/DagOnStar/dagonstar}) keeps
its own example families: taskflow and dataflow examples, Docker, Slurm,
cloud and Globus dataflows, transversal and meta-workflows, the HiPES
tutorial, environmental applications, and a 36-lesson tutorial
(\texttt{docs/tutorial/}, \texttt{examples/tutorial/}). They are a natural
next step once the course examples are familiar.

\section{Recommended Reading and Verification Order}

The documentation suggests beginning with examples that are verifiable on an
ordinary developer machine, then progressing to examples that require
site-dependent services:

\paragraph{Stage 1 --- Local, no external service.}
\begin{enumerate}[leftmargin=1.6em]
  \item Taskflow (\Cref{ch:lessons-01-02});
  \item Batch dataflow (\Cref{ch:lessons-03-04});
  \item Checkpointing (\Cref{ch:lessons-05-07}, Appendix~D);
  \item Docker, if a Docker daemon is available locally
    (\Cref{ch:backends}).
\end{enumerate}

\paragraph{Stage 2 --- Site-dependent.}
\begin{enumerate}[leftmargin=1.6em]
  \item Slurm (\Cref{ch:backends});
  \item Remote batch (\Cref{ch:backends});
  \item Cloud;
  \item Globus;
  \item Environmental application (\Cref{ch:case-study}).
\end{enumerate}

\begin{practicebox}
This staged strategy is itself worth adopting as a general engineering
practice when developing a \emph{new} DAGonStar workflow, not only when
reading the example catalog: build and verify the local-only skeleton of a
pipeline first --- using mock services where the LLM and Web lessons already
demonstrate the pattern (\Cref{ch:llm-tasks,ch:web-tasks}) --- before
introducing the site-specific services (a real Slurm cluster, a real cloud
account, a real LLM provider) that the production version will ultimately
require.
\end{practicebox}

\section{The Consistency Rule}

The examples directory includes legacy research material alongside actively
maintained tutorials. The documentation distinguishes four categories
explicitly, and this book preserves the same distinction:

\begin{itemize}[leftmargin=1.4em]
  \item locally verifiable examples;
  \item examples requiring external software;
  \item examples requiring remote services or site-specific credentials;
  \item historical scripts retained for reproducibility context.
\end{itemize}

\begin{cautionbox}
A reader extending these examples toward production use should explicitly
re-classify any adapted example into one of these four categories and record
that classification alongside the adapted workflow --- precisely the kind of
operational metadata that keeps a scientific pipeline auditable as it
evolves away from its tutorial origin.
\end{cautionbox}

\section{The Pedagogical Function of an Examples Catalog}

It is worth situating the examples catalog's organizational philosophy
within a broader observation about how research software documentation
tends to fail, and how a catalog of the kind \Cref{ch:examples-catalog}
surveys is designed specifically to avoid that failure mode.

A common pattern in research software projects is for examples to
accumulate organically, over years, as individual researchers contribute
scripts addressing their own immediate needs, with no consistent
classification of what each example assumes, requires, or guarantees. The
result, familiar to anyone who has inherited a research codebase, is a
directory of scripts of wildly varying reliability, with no way to tell,
short of reading each one in full, which examples will run unmodified on a
fresh checkout and which silently assume a specific cluster, credential set,
or dataset the original author had access to but never documented. The
four-category consistency rule \Cref{ch:examples-catalog} reproduced ---
locally verifiable; requiring external software; requiring remote
services or credentials; historical, retained for context --- exists
precisely to prevent this accumulation from becoming opaque: every example
is required to declare, explicitly, which category it belongs to, so that a
reader can decide, before investing time, whether a given example is even
attemptable in their current environment.

\begin{practicebox}
This four-category discipline is worth adopting independently of DAGonStar,
as a general practice for any research group maintaining a shared example or
script collection: require every contributed example to declare its
verification category at the point of contribution, rather than leaving
later readers to discover the category empirically, through failed attempts
to run code whose external requirements were never documented.
\end{practicebox}

\section{Mock Services as a Verification Strategy}

The staged verification strategy of \Cref{ch:examples-catalog} --- local
examples first, site-dependent examples second --- depends critically on a
technique this book has used repeatedly without naming it explicitly until
now: the \emph{mock service}. Lectures~25--30's locally verifiable web, LLM and FaaS
examples (\Cref{ch:llm-tasks,ch:web-tasks}) all achieve local
verifiability by substituting a small, purpose-built HTTP server --- bound
only to \texttt{127.0.0.1}, requiring no credential, implementing only the
minimal subset of a real provider's protocol surface needed for the
example --- for the real, external service a production workflow would
ultimately target.

This technique deserves recognition as a general-purpose strategy, not
merely an implementation convenience specific to the tutorial. A mock
service verifies the \emph{mechanism} --- dependency inference, staging,
request construction, response capture --- entirely independently of the
\emph{content} a real provider would return, which is exactly the
separation of concerns a researcher needs when developing or debugging a
new \texttt{LLM}- or \texttt{WEB}-task-based workflow stage: mechanism bugs
(a malformed request, a missing dependency, an incorrectly staged file) can
be found and fixed against the mock, with zero cost and zero external
dependency, before ever spending a real API quota or risking a real,
possibly non-idempotent, side effect against a production endpoint.

\begin{example}
A researcher developing a new \texttt{WEB}-task-based integration with an
unfamiliar third-party API can profitably begin by writing a mock server
that returns a fixed, known-good response matching the target API's
documented response schema, developing and debugging the DAGonStar task
specification entirely against this mock, and only substituting the real
endpoint URL --- with no other change to the task specification, by the
provider-agnostic design principle already discussed in
\Cref{ch:llm-tasks} --- once the mechanism is confirmed correct. This
mirrors exactly the relationship between Lecture~27's \texttt{local\_mock}
provider and Lecture~27's documented path to a real provider
(\Cref{ch:llm-tasks}, Step~6).
\end{example}

\section{Reading an Unfamiliar Example Critically}

Building on the recommended reading order of \Cref{ch:examples-catalog}, it
is worth giving a concrete procedure for how a reader should approach any
\emph{individual} example in the catalog, beyond simply running it, since
the goal of reading an example is generally to extract a reusable pattern
for one's own workflow, not merely to confirm that the example itself
executes successfully.

\begin{enumerate}[leftmargin=1.8em]
  \item \textbf{Identify the verification category first}
    (\Cref{ch:examples-catalog}), to calibrate what can be assumed about
    the example's reliability and portability before reading further.
  \item \textbf{Identify which task types and staging modes the example
    exercises}, using the vocabulary developed across Parts~I--IV of
    this book, to situate the example relative to material already covered.
  \item \textbf{Identify the dependency structure}, ideally by sketching the
    DAG the example constructs, before reading the task bodies in detail
    --- the structural skeleton is usually more transferable to a reader's
    own problem than the specific commands inside any one task.
  \item \textbf{Identify what the example deliberately omits}, particularly
    for examples in the ``requiring external software'' or ``requiring
    remote services'' categories, where production-readiness concerns
    (\Cref{ch:case-study}'s operational-deployment discussion is directly
    relevant here) are frequently out of scope by design, for the sake of
    pedagogical focus.
\end{enumerate}

\section{Chapter Summary}

The examples catalog is organized around the same staged, local-first
verification philosophy this book has followed since \Cref{ch:lessons-01-02}:
prove a mechanism with no external dependency, then extend it toward the
real service it is meant to integrate with. This chapter additionally
situated the catalog's four-category consistency rule within a broader
critique of how research-software examples typically accumulate without
documentation, named and generalized the mock-service verification
technique already used implicitly by the LLM and Web task lessons, and gave
a concrete, four-step procedure for critically reading any unfamiliar
example in the catalog for transferable patterns rather than merely
confirming its successful execution. \Cref{ch:synthesis} now closes the
book by assembling every mechanism surveyed across Parts~I--IV into a
coherent picture of an operational, reproducible, AI-augmented
forecasting chain.

\begin{exercise}
Apply the four-step reading procedure of this chapter to the HiPES tutorial
examples family of the upstream repository, based
only on its one-line catalog description (``a staged scientific tutorial
from configuration to map rendering''). What predictions can you make about
its likely task-type composition and dependency structure before consulting
the source material directly, and what aspects of the four-step procedure
necessarily require reading the actual source rather than the catalog
description alone?
\end{exercise}
